% ============================================================
%  chapters/ch03-methodology.tex
% ============================================================

\chapter{Methodology}
\label{ch:methodology}

\section{Overview of Proposed Framework}

% Describe your methodology.

\section{Dataset Description}
\label{sec:datasets}

\section{Proposed Architecture}
\label{sec:arch}

% Example figure:
%
% \begin{figure}[H]
%   \centering
%   \includegraphics[width=0.85\textwidth]{images/architecture.png}
%   \caption{Overview of the proposed architecture.}
%   \label{fig:arch}
% \end{figure}

\section{Loss Function}
\label{sec:loss}

\section{Training Details}
\label{sec:training}

\begin{table}[H]
  \centering
  \caption{Training hyperparameters}
  \label{tab:hyperparams}
  \begin{tabular}{ll}
    \toprule
    \textbf{Parameter} & \textbf{Value} \\
    \midrule
    Optimiser     & Adam \\
    Learning rate & $1 \times 10^{-4}$ \\
    Batch size    & 16 \\
    Epochs        & 100 \\
    \bottomrule
  \end{tabular}
\end{table}

\section{Evaluation Metrics}
\label{sec:metrics}

The following metrics are used to evaluate performance.

\begin{description}
  \item[PSNR] Peak Signal-to-Noise Ratio measures pixel-level fidelity.
  \item[SSIM] Structural Similarity Index captures perceptual similarity.
  \item[EPI] Edge Preservation Index quantifies retention of high-frequency
    structural information.
\end{description}
